\section{Method}\label{sec:method}

\subsection{Model and contrastive-correctness adapter}\label{method:model}

All results are on Gemma~4 31B IT, run in bf16 --- the Hugging Face checkpoint
\code{google/\allowbreak gemma-\allowbreak 4-\allowbreak 31B-\allowbreak it}, a
\code{Gemma4ForConditionalGeneration} with 60 decoder layers, $d_{\mathrm{model}}$
5376, 32 query heads and head dimension 256, evaluated under
\code{transformers}~5.8.0 and \code{lm-eval}~0.4.12~\citep{LMEvalHarness}.

\paragraph{Which snapshot, and why it is two.} Our harness resolves the model by
\code{revision=main} and lm-eval leaves \code{model\_sha} empty, so no result file
carries an explicit pin. The snapshot is nonetheless recoverable for most runs,
because the harness records the model path verbatim and on our machines that path is a
Hugging Face cache directory containing the snapshot hash; we extract it from every
artifact and release the resulting map (\code{scripts/\allowbreak model\_\allowbreak provenance.\allowbreak py},
\code{runs/\allowbreak provenance/}). The exception is the rented-GPU campaigns, which
read a pre-staged copy under a flat path (\code{/workspace/\allowbreak models/}) that
carries no hash --- the canonical June trio among them. Those we date rather than
resolve, which is weaker evidence and we mark it as such. Upstream published a revised
checkpoint on 2026-07-20, so \code{main} moved mid-project and our runs straddle it.
The June and early-July campaigns --- the arbiter, the inducibility control, the
recovery spectrum, the on-axis reference, and, by date rather than by hash, the
headline trio (run 2026-06-10, four weeks before the newer template existed) --- ran on
\code{3548789868c5356dbf307c98e6f609007b82b3eb}; the late-July campaigns --- the
full-split scoring replication of \cref{disc:status}, the anti-axis arms and the
dense scale sweep --- ran on \code{842da3794eaa0b77d5f08bae87a17459d91ff475}, which
is \code{main} at the time of writing. We diffed the two snapshots file by file:
\code{config.json} and \code{tokenizer.json} are byte-identical, so the architecture
and the vocabulary did not change, while \code{chat\_template.jinja} and
\code{tokenizer\_config.json} differ --- the newer template is dated 2026-07-09 and
its own header describes fixes to tool-calling loops and turn closures. Because every
protocol here is template-bound (\cref{results:headline}), we flag this rather than
bury it. It cuts in the reassuring direction: the headline scoring lift is observed
under \emph{both} templates --- the $n{=}400$ campaign under the older one and the
$n{=}1172$ full-split replication ($+36.26\pp$) under the newer --- so the central
effect survives the change rather than depending on it. The remaining arms were each
run under one snapshot only, and we do not claim more than that. We accordingly treat
snapshot-level reproducibility as a limitation of these runs, and the released
harness now stamps the resolved snapshot and a hash of the applied chat template into
every result file.

The cross-family replication
(\cref{sec:limitations}) uses \code{Qwen/\allowbreak Qwen2.5-\allowbreak 14B-\allowbreak Instruct}
at revision \code{cf98f3b3\allowbreak bbb457ad\allowbreak 9e2bb7ba\allowbreak f9a0125b\allowbreak 6b88caa8}
(\code{Qwen2ForCausalLM}, 48 layers, $d_{\mathrm{model}}$ 5120); that repository has
been unmodified since September 2024, well before these runs, so \code{main} and the
pinned revision coincide. The discrimination
adapter (referred to as V7-mc) is a contrastive-correctness PEFT in the LoFiT/ITI
family: it learns
additive per-head offset vectors on a sparse set of attention heads from contrastive
pairs of correct and incorrect answers under a multiple-choice discrimination
objective. Concretely, each pair is (question, correct option, first-listed
incorrect option), built from the TruthfulQA mc1 validation split (seed-0
permutation, 80\% train side) and the native ARC-Challenge train split; training
minimizes a margin loss on the summed continuation log-probability under the chat
template, $-\log p(\mathrm{correct}) + \gamma\cdot\log p(\mathrm{wrong})$ with
$\gamma = 0.5$, for 800 steps, with the deployed checkpoint selected by a
validation-accuracy monitor. Two construction details a reader of the code will
find, stated here rather than left to be discovered. First, because the HF
TruthfulQA mc1 release lists the gold answer at index 0, ``first-listed incorrect
option'' means every TQA pair contrasts the dataset's first two options --- a fixed
positional split rather than a sample over distractors; this touches no mechanism
claim (those are anchored on ARC, where option order carries no gold information)
and is consistent with --- plausibly contributes to --- TruthfulQA's answer-prior
character (\cref{disc:heterogeneity}). Second, the checkpoint-selection monitor was,
due to dataset ordering in the validation slice, computed on TruthfulQA items only:
checkpoint selection never saw ARC signal, so the ARC lift is not
checkpoint-selected. The causal analyses use the 48 heads carrying the adapter's
offset (selection criterion in \cref{sec:definitions}). Those 48 (layer, head)
targets are unique (no aliasing) and the offset modulates a sparse,
\emph{non-contiguous} set of 12 layers --- [30, 31, 32, 33, 34, 36, 37, 38, 39, 40,
49, 55] (deterministic across runs; note the gap at layer 35) --- so references
below to a ``layer ${\sim}$30--55 span'' denote this discrete set's bracket, not a
contiguous band.

\subsection{Evaluation protocol: scoring-face vs generation-face}
\label{method:protocol}

We distinguish two protocols and report absolutes only with their protocol attached.
\textbf{Scoring-face} is cold multiple-choice scoring: log-likelihood ranking of
options with the chat template applied and no generation (lm-eval canonical,
\code{-\allowbreak -\allowbreak apply-\allowbreak chat-\allowbreak template}, $n{=}400$, acc\_norm), with identical flags across
arms (one documented exception, the random-offset control, in
\cref{method:controls}). The $n{=}400$ sample is the first 400 items of each task's evaluation split in
native order (lm-eval \code{-\allowbreak -\allowbreak limit 400}), not a random draw --- identical across
arms by construction, and not doing selection work: the R4 full-split run confirms
every headline delta within the $n{=}400$ confidence interval
(\cref{results:headline}, \cref{sec:limitations}). This is deliberate --- it yields
apples-to-apples deltas against the adapter --- and it is why our base absolutes
(e.g.\ ARC ${\approx}0.50$) are scoring-face floors and must not be compared to a
third-party-reported ${\sim}91\%$ chain-of-thought~\citep{ChainOfThought} figure (the protocol caveat in
the \cref{tab:headline} caption). The same scoring-face protocol --- identical
log-likelihood option ranking, chat template, and lm-eval canonical flags --- is
also applied out-of-task to GPQA~\citep{GPQA}, where the cold-MC lift is itself ${\sim}0$ (base
$0.4747 \rightarrow$ adapter $0.4697$, $-0.5\pp$); this cold-MC GPQA measurement is
distinct from the generation-face GPQA evaluation described below.
\textbf{Generation-face} evaluates the produced text: gsm8k~\citep{GSM8K} under
\code{generate\_\allowbreak until} with exact-match; a per-item arbiter that runs base
generative chain-of-thought (thinking-on, greedy, $k{=}1$) aligned to the cold-MC
\code{doc\_\allowbreak id}s, on the items the adapter fixes; and GPQA under thinking-on greedy
decoding with the generation budget raised to 8192 tokens (at 2048 the base
truncated 96\% of generations, producing a spurious 0.15). The per-item arbiter is
run across multiple anchors --- ARC, GPQA, HellaSwag, and Winogrande --- to test
whether the fixed items are already CoT-reachable. One caveat is load-bearing here:
HellaSwag and Winogrande are not native letter-MCQ tasks, so to run the generative
arbiter on them we reformulate each item as a lettered MCQ (HellaSwag: the four
continuations $\rightarrow$ A/B/C/D; Winogrande: the two options $\rightarrow$ A/B).
That reformulation puts the arbiter in a \emph{different} format from the cold-MC
scorer (which scores option continuations without letters), so on
HellaSwag/Winogrande the arbiter anchors \emph{knowledge} in a letter-MCQ format
rather than replicating the cold-MC channel --- the same format gap that ARC and
GPQA already carry (where the arbiter is nonetheless informative). We therefore read
the arbiter as a knowledge-accessibility probe, not as a channel-faithful readout.

\subsection{Functional correctness direction and on-axis references}
\label{method:directions}

Per attention head we estimate a functional correctness direction $\vinf$ (the
mass-mean difference between correct and incorrect inference-form activations). The
estimation corpus is the inference-family contrastive pairs (ARC-Challenge native
train split plus HellaSwag), so on ARC --- where every on-axis recovery fraction is
anchored --- the items that estimate the direction and the test items that measure
recovery are disjoint by construction, the same train/test boundary the adapter
itself respects (\cref{sec:limitations}). To make the on-axis causal test as strong
as possible, we also construct $\Wknow$, a whitened Fisher-LDA correctness direction
--- the strongest on-axis reference we can build locally, and genuinely distinct
from the mass-mean estimate (cos with $\vinf$ 0.379). ``Strongest'' is a
construction criterion, not an outcome criterion: the whitened Fisher-LDA direction
is the optimal linear correctness discriminant under the class covariance, which the
raw mass-mean ignores --- it is not chosen for recovering the least lift, and
\cref{disc:geometry} reports both references' recovery side by side. Because
optimality under the class covariance is a population statement and the covariance
here is estimated, we also check the two references \emph{as readers} rather than
arguing from construction alone: on the base-wrong cold-MC population both separate
gold from distractors at or above the level of a probe fitted on that same space
($\Wknow$ $0.9613$, $\vinf$ $0.9646$, fitted probe $0.9545$), and paired over the
same items the gap between them is $-0.0034$ with a CI95 of $[-0.0135, +0.0067]$ ---
so ``strongest'' is not contradicted by outcome, and neither is it established by
it (\cref{disc:geometry}).

\textbf{Scope of ``off-axis''.} Both references are \emph{one direction per head}, not
a single global axis: the geometry is a family of per-head linear correctness axes, one
at each intervened head. Everything we call off-axis is therefore off the measured
linear correctness axis at each intervened head, which is the strongest such direction
we can build there but not the whole geometry of correctness: correctness information has been reported to occupy a
low-dimensional subspace of two to eight directions rather than one, with a causal
rank bounded at five under distributed alignment search
\citep{ConfidenceManifold}. Our negative arms show that the learned offset is not
carried by the best 1-D discriminant; they do not show that it lies outside every
correctness-bearing subspace. Testing $k$-dimensional subspaces with norm-matched
parallel and residual re-injection is the natural follow-up and we do not claim its
result here. The on-axis
causal control substitutes an offset built along $\vinf$ (and along $\Wknow$), unit
direction scaled to the per-head norm of the trained offset across all 48 heads, so
that \emph{direction} is the only variable, gated by a sign-safe three-zone control
gate with a null-task check (the gate, and the Gate-A reproduction check it
presupposes, are defined in \cref{sec:definitions}).

\subsection{Net-effect geometry and the inducibility control}\label{method:neteffect}

To compare methods fairly in activation space (rather than via the weight-space
offset, which we show is a misleading proxy), we measure each method's \emph{net
effect on activations} and decompose it against $\vinf$, reporting the
off-axis/on-axis split as a ratio relative to a random direction.

The inducibility question --- what fraction of the contrastive lift a
\emph{non-contrastive} objective recovers --- is asked with two controls, because the
answer depends on how much capacity the control is given.

\textbf{(i) Same-parameterization control (the objective, isolated).} The identical
LoFiT apparatus as the adapter: the same 48 head-pairs from the same selection file,
the same $12{,}336$ trainable parameters ($\alpha_h \in \mathbb{R}$,
$\theta_h \in \mathbb{R}^{256}$ per head), the same corpus, optimizer steps, batch
size, learning rate, sequence length, train/validation split and best-checkpoint
rule. The \emph{only} change is the loss: ordinary cross-entropy on the correct
continuation, $-\log p(\text{correct})$, with no term pushing down the incorrect
option. Three seeds vary the batch-visit order with the data split held fixed, so the
three runs see the same examples. Because the intervention space is held exactly
fixed, this control attributes the difference to the \emph{objective} and nothing
else. That equality is recorded inside the released checkpoints rather than inferred:
each offset file carries the training \code{config}, whose \code{loss\_type} reads
\code{direct} for the adapter and \code{plain\_ce} for the three control seeds, with
the same head-selection file, the same $48$ (layer, head) pairs in the same order, and
the same \code{top\_k}, step count, learning rate and corpus; the $12{,}336$ is counted
on the tensors themselves ($\theta \in \mathbb{R}^{48 \times 256}$ plus
$\alpha \in \mathbb{R}^{48}$), not inferred from file size, and the released training
log states it independently. One scope note: the adapter's \code{config} predates the
field that records the training seed, so the equality is verified over the fields both
sides record --- which are all of those that fix the parameterization --- and not over
that one.

One asymmetry in that control deserves stating, since it touches the arm the
attribution rests on. The \emph{objective} is genuinely non-contrastive --- the
gradient carries no term against the incorrect option --- but the
\emph{checkpoint-selection rule}, which we held identical to the adapter's precisely
to keep the apparatus fixed, is not: validation tracks whether the correct option
outscores the incorrect one, so the distractor enters the choice of which step to
keep even though it never enters the loss. So ``only the loss changed'' is exact for
the training signal and not for model selection. Two things bound the consequence.
The rule only bit in one of the three seeds --- in seeds 1 and 2 the selected
checkpoint \emph{is} the final one, and only in seed 0 does the best-by-validation
step (600) differ from the last (800). And in that seed the rule selects the
\emph{worse} ARC checkpoint, which is why the fixed-final variant, which consults no
validation signal at all, reports a \emph{higher} recovered fraction
($39.7\% \pm 3.1$) than the best-checkpoint figure we quote ($36.2\% \pm 5.0$). The
contamination therefore deflates this control rather than inflating it, and the
residual we attribute to the contrastive objective is, if anything, conservative. We
quote the best-checkpoint number as the headline and carry the ${\sim}36$--$40\%$
range throughout.

\textbf{(ii) Higher-capacity control (ordinary SFT at scale).} A non-contrastive
Align-LoRA: ordinary supervised fine-tuning on a data- and step-matched corpus,
swept over LoRA rank (r8/32/64/256) and scored through the same three-zone control
gate. This arm is \emph{not} capacity-matched --- LoRA on the text decoder's
q/k/v/o projections across 60 layers carries ${\sim}22$M trainable parameters at r8
and ${\sim}1.44$B at r512, i.e.\ $1{,}800\times$ to $117{,}000\times$ the adapter's
$12{,}336$ --- and we report it as such. The recovery fraction we quote
(${\sim}62\%$) is the r8--r256 plateau, i.e.\ $1{,}800\times$--$58{,}000\times$; the
r512 point is a separate best-effort flank and is always reported as one. The mismatch is conservative for the
inducibility claim (a control with far more capacity that still recovers only part of
the lift), but it is the reason the same-parameterization arm exists: only (i)
licenses reading the residual as contrastive-specific \emph{at equal capacity}.

Both recovery fractions, the matched corpus, and the gate are defined in
\cref{sec:definitions}.

\subsection{Additional controls}\label{method:controls}

We further run: a norm-matched random offset (three seeds) to test whether any
perturbation of the right size reproduces the lift --- this arm is the one exception
to the identical-flags statement of \cref{method:protocol}, having been scored with a
2048-token context window against 4096 for the canonical trio and the other spectrum
arms; task and \code{-\allowbreak -\allowbreak limit} match, and no truncation is
reachable at either setting because the longest prompt in this protocol is
${\approx}228$ tokens (\cref{sec:limitations}), but we state it rather than let
``identical flags'' cover it; a same-scaffold PMI /
decision-conditional decomposition and a choices-only (question-removed) control to
separate question-conditional gains from answer priors
\citep{SurfaceFormCompetition}; layer-band patching (restricting the offset to layer
windows) to test for a layer locus; and a within-format difficulty stratification on
ARC --- ARC-Easy vs ARC-Challenge under the identical cold-MC scoring protocol
(acc\_norm), stratified by base-wrong items --- which holds format and distribution
fixed and varies only difficulty, isolating the contribution of difficulty from that
of distribution distance.

\subsection{Formal definitions}\label{sec:definitions}

This subsection makes the operational terms used throughout precise --- each
grounded in the analysis code. The maps from every headline number to the run and
the artifact that produced it, so no two item populations are silently conflated,
are in \cref{app:provenance}.

\paragraph{Capability.}
Throughout, \emph{capability} is used in a readout-independent sense: accuracy the
base attains under \emph{any} readout we measured, at comparable or lower cost. A
reader who reserves the term for a fixed readout will read our claim as being about
the cold-MC channel specifically --- and on that reading the adapter does improve it,
which is the dissociation this paper is about rather than an objection to it.

\paragraph{The trained offset $\thetamc$.}
$\thetamc$ denotes the adapter's learned per-head additive offset --- the aggregate
write direction across the 48 heads (\cref{method:model}). The geometric and causal
decompositions of \cref{disc:geometry} ($\offpar$/$\offperp$, whitened cosines)
operate on it.

\paragraph{Recovery fraction.}
For each anchor task the recovery fraction of an arm is
$\mathrm{frac} = (\mathrm{arm} - \mathrm{base}) / (\mathrm{v7mc} - \mathrm{base})$,
the fraction of the contrastive lift the arm reproduces
(\code{scripts/\allowbreak analyze\_\allowbreak r1\_\allowbreak align\_\allowbreak frac.\allowbreak py:6,99}; returns NaN when
$|\mathrm{v7mc} - \mathrm{base}| < 10^{-9}$). The per-task headline metric is fixed
in code (\code{TASK\_\allowbreak METRIC}, \code{:36-\allowbreak -\allowbreak 41}): \textbf{acc\_norm} for ARC-Challenge~\citep{ARC}
and HellaSwag~\citep{HellaSwag}, plain \textbf{acc} for Winogrande~\citep{WinoGrande}
and TruthfulQA-mc1~\citep{TruthfulQA} (the latter two
define no length-normalized variant). Every ``fraction of the lift'' in this paper
is this quantity at these metrics.

\paragraph{Gate-A (canonical reproduction).}
Before any causal arm is trusted, \code{scripts/\allowbreak check\_\allowbreak canonical\_\allowbreak repro.\allowbreak py} re-runs
the canonical R1 baselines and aborts unless they reproduce: in \emph{reproduce}
mode $|\mathrm{measured} - \mathrm{expected}| \leq 0.05$ for both the base
(ARC-Challenge acc\_norm 0.505) and the adapter (0.855) (\code{:69,87-\allowbreak -\allowbreak 91}); in
\emph{null} mode the zero-offset wrapper ($\alpha{=}\theta{=}0$) must be
behaviorally identical to the base, $|\mathrm{null} - \mathrm{base}| \leq 0.02$
(\code{:16}), certifying the lift is not a wrapper artifact. A FAIL exits non-zero
and blocks the dependent batch. The Gate-A base values are the exact reference the
on-axis arms are compared against (e.g.\ $\offpar$ recovering to 0.505 = ``0\%'',
\cref{disc:geometry}).

\paragraph{Three-zone control gate.}
The inducibility fraction is read through a sign-safe three-zone gate
(\code{scripts/\allowbreak check\_\allowbreak canonical\_\allowbreak repro.\allowbreak py:70-\allowbreak -\allowbreak 150}) that keeps numerator and
denominator harness-internal (the native HFLM for the control's own effect, the V7
wrapper for the contrastive lift) so no cross-harness bias enters the ratio. Two
sign-safety checks precede any fraction: a \textbf{setup} gate requires the
contrastive lift to be real and positive ($\mathrm{v7mc} - \mathrm{base} >
\mathrm{min\_effect}$, min\_effect 0.01) and a \textbf{convergence} gate requires
the control's own effect to be non-trivial and positive
($\mathrm{arm} - \mathrm{base} \geq 0.01$), so a no-op or harmful control cannot be
scored. The surviving fraction is binned at $\mathrm{frac\_lo} = 1/3$ and
$\mathrm{frac\_hi} = 2/3$ into three zones (\code{:71,73,139-\allowbreak -\allowbreak 150}); the
${\sim}62\%$ ARC recovery lands in the middle zone, which is what we report as
``zone C'' against these thresholds. The thresholds bin the \emph{point} estimate, so
the zone assignment does not depend on which interval accompanies it (paired
bootstrap $[51\%, 73\%]$, delta-method $[45\%, 85\%]$; \cref{disc:inducibility}).
A separate null-task check (MMLU~\citep{MMLU}) guards that a causal arm
does not damage an unrelated task (null MMLU $+0.07\pp$ at PASS,
\cref{tab:onaxis}).

\paragraph{Whitened cosine.}
Geometric alignments (e.g.\ $\cos(\thetamc, \vinf) = -0.003$,
$\cos(\vinf, \vret) = +0.283$) are Mahalanobis cosines under the per-head pooled
class-activation covariance, in the 256-dim per-head activation space. The canonical
headline values come from an eigendecomposition whitening
(\code{scripts/\allowbreak probe\_\allowbreak oq1\_\allowbreak functional\_\allowbreak angle.\allowbreak py::\_\allowbreak whitening\_\allowbreak from\_\allowbreak cov}, reused
verbatim by \code{characterize\_\allowbreak theta\_\allowbreak mc.\allowbreak py}):
$W = V\cdot\mathrm{diag}(\lambda^{-1/2})\cdot V^{\top}$ from the eigh of
$\mathrm{cov\_common} = \tfrac{1}{2}(\mathrm{cov\_inf} + \mathrm{cov\_ret})$,
eigenvalues floor-clamped at $10^{-6}$ --- a floor that never binds in the actual
computation (zero clamped eigenvalues; median condition number ${\sim}894$). Each
per-head covariance is estimated from $N = 2400$ activation samples (1200 correct +
1200 incorrect) against $d = 256$ dimensions, an $N/d$ ratio of ${\sim}9.4$. A
related additive-ridge formulation,
$M = (\mathrm{cov} + 10^{-2}\cdot\mathrm{mean}(\mathrm{diag\,cov})\cdot I)^{-1}$,
appears in a separate probe-direction sanity script and (as a shrinkage on the
covariance) in the $\offpar$/$\offperp$ construction; it is not the source of the
headline cosines. Robustness: sweeping that ridge factor over two orders of
magnitude ($\{10^{-3} \ldots 10^{-1}\}$) moves $\cos(\thetamc, \vinf)$ only within
$[-0.0032, -0.0021]$ (always inside the null band, sign stable --- stronger
regularization pushes it \emph{toward} zero) and $\cos(\vinf, \vret)$ within
$[+0.284, +0.325]$ (sign and magnitude stable), so neither headline geometric fact
is an artifact of the regularization choice
(\code{runs/\allowbreak oq1\_\allowbreak functional\_\allowbreak axis/\allowbreak ridge\_\allowbreak sensitivity.\allowbreak json}).

\paragraph{The 48 heads.}
The causal analyses use the 48 (layer, head) cells carrying the adapter's offset.
These are the \textbf{top-K = 48} heads ranked by per-head probe accuracy
(\code{best\_\allowbreak acc} = the max of the Fisher- and Ridge-CV probes), a pure ranking
slice with \textbf{no accuracy-threshold gate}
(\code{scripts/\allowbreak probe\_\allowbreak v7\_\allowbreak lofit\_\allowbreak head\_\allowbreak selection.\allowbreak py:1029-\allowbreak -\allowbreak 1032}); the 0.60 accuracy
threshold in that script only annotates the report, it does not select heads
($K = 48$ follows the LoFiT recommendation).

\paragraph{Align-LoRA corpus and budget.}
The inducibility control is a non-contrastive LoRA fine-tune under a plain
causal-LM cross-entropy loss (\code{scripts/\allowbreak train\_\allowbreak align\_\allowbreak lora.\allowbreak py}), on a corpus
\textbf{matched to the V7-mc training data}: the matched builder
(\code{scripts/\allowbreak build\_\allowbreak align\_\allowbreak control\_\allowbreak corpus.\allowbreak py}) imports only the TruthfulQA and
ARC splits used to train V7-mc, byte-identically (TQA = the first 80\% of the
validation split under a seed-0 permutation; ARC = the native ARC-Challenge train
split), and the control is swept over LoRA rank $r \in \{8, 32, 64, 256\}$ at 2
epochs, effective batch 8 (batch 1 $\times$ grad-accum 8), max sequence 1024.
\textbf{``Budget-matched'' means the same training examples and the same optimizer
steps as the contrastive adapter --- not the same FLOPs} (the two objectives differ
in per-step cost). ARC is the clean anchor; TruthfulQA is excluded from the fraction
because its matched corpus overlaps its evaluation (\cref{disc:inducibility}).
